% GENERATED FILE. Edit canonical lessons, then run npm run generate:books.
% canonical-source-hash: fnv1a64:f5549ec87595f83a
% canonical-lessons: ZH-C422-daohang, ZH-C422-dongbei, ZH-W422-bei-north, ZH-C422-xibei, ZH-C422-beishang, ZH-C422-kuaisu

\chapter{Going North}
\label{ch:zh-going-north}
\hlchaptermodality{422}

\begin{chapteropening}
\textbf{By the end of this chapter:} \emph{I can use a satnav, name the northeast and the northwest, say we are heading north and take a fast bus.}

Complete the last lesson of chapter 422: I can use a satnav, name the northeast and the northwest, say we are heading north and take a fast bus.
\end{chapteropening}

\section[dǎoháng]{dǎoháng — to navigate by satnav}

\label{lesson:ZH-C422-daohang}



\begin{quote}
Before the new one: say the Mandarin for a toll gate, then the Mandarin for to break the speed limit.
\end{quote}

\subsection*{You'll want to know: dǎoháng}
\emph{dǎoháng} — \textquotedblleft{}to navigate by satnav\textquotedblright{}.

\begin{quote}
\emph{wǒ yòng shǒujī dǎoháng zhǎodào le nà jiā fàndiàn} — I found the way to that restaurant with my phone.
\end{quote}

Its characters come later; for now it is a word for the ear.

\begin{cousinweb}
\emph{dǎoháng} is built from \emph{dǎo} 'to lead' and \emph{háng} 'to sail', both inherited from Middle Chinese, \emph{dawH} and \emph{hang} (Guangyun, in Baxter's transcription), which glosses the first 'to lead' and the second 'a boat'. \emph{dawH} would regularly give dào: the falling-rising tone of dǎo is irregular.
\end{cousinweb}

\begin{grammarlens}[title={What you've built}]
One more word for getting around.
\end{grammarlens}

\subsection*{Your turn}
\begin{itemize}
\raggedright
  \item \emph{Say it:} \emph{dǎoháng}
  \item \emph{Say it:} \emph{dǎoháng}, once more
  \item \emph{Say it:} \emph{wǒ yòng shǒujī dǎoháng zhǎodào le nà jiā fàndiàn} — I found the way to that restaurant with my phone.
  \item \emph{Recall:} say the Mandarin for a toll gate, then the Mandarin for to break the speed limit, then say \emph{dǎoháng} again
\end{itemize}

\begin{tcolorbox}[breakable,colback=teal!4,colframe=teal!35!black,title={Before you move on}]
What does dǎoháng mean? (\textquotedblleft{}To navigate by satnav\textquotedblright{}.) Say it once more.
\end{tcolorbox}
\section[\texorpdfstring{\zh{东北} (dōngběi)}{东北 (dōngběi)}]{\zh{东北} (dōngběi) — northeast; the Northeast (of China)}

\label{lesson:ZH-C422-dongbei}



\begin{quote}
Before the new one: say the Mandarin for to break the speed limit, then the Mandarin for to navigate by satnav.
\end{quote}

\subsection*{You'll want to know: \zh{东北}}
\textbf{\zh{东北}} — \emph{dōngběi} — \textquotedblleft{}northeast; the Northeast (of China)\textquotedblright{}.

\begin{quote}
\textbf{\zh{火车进入了东北}} \emph{huǒchē jìnrù le dōngběi} — The train has come into the northeast of China.
\end{quote}

The new character, in pieces: \zh{北} — five strokes in two halves: on the left a vertical, a short bar into it from the left and a bar rising to the right at its foot; on the right a stroke falling to the left, then a vertical that bends right along the foot and hooks up. The next lesson puts it in your hand.

\begin{cousinweb}
\textbf{\zh{东北}} (\emph{dōngběi}) is built from \zh{东} 'east' and \zh{北} 'north', both inherited from Middle Chinese, \emph{tuwng} and \emph{pok} (Guangyun, in Baxter's transcription), which glosses the second 'north' and 'to flee'.
\end{cousinweb}

\begin{grammarlens}[title={What you've built}]
One more word for getting around.
\end{grammarlens}

\subsection*{Your turn}
\begin{itemize}
\raggedright
  \item \emph{Say it:} \emph{dōngběi}
  \item \emph{Say it:} \emph{dōngběi}, once more
  \item \emph{Say it:} \emph{huǒchē jìnrù le dōngběi} — The train has come into the northeast of China.
  \item \emph{Recall:} say the Mandarin for to break the speed limit, then the Mandarin for to navigate by satnav, then say \emph{dōngběi} again
\end{itemize}

\begin{tcolorbox}[breakable,colback=teal!4,colframe=teal!35!black,title={Before you move on}]
What does \zh{东北} mean? (\textquotedblleft{}Northeast; the Northeast (of China)\textquotedblright{}.) Say it once more.
\end{tcolorbox}
\section[\texorpdfstring{\zh{北} (běi)}{北 (běi)}]{\zh{北} — two halves back to back}

\label{lesson:ZH-W422-bei-north}



\begin{quote}
The word first, then the character inside it.

\begin{itemize}
\raggedright
  \item \emph{Recall:} say \emph{dōngběi} — \textbf{R1}, one lesson back, before it decays
\end{itemize}

\begin{quote}
\zh{比} \zh{东}
\end{quote}

\emph{Read it:} each one aloud, with nothing said first — \emph{bǐ}, \emph{dōng}
\end{quote}

\subsection*{Script — two halves back to back}
\hlblockfigure{\detokenize{figures/ZH-W422-bei-north-filmstrip.pdf}}{How \zh{北} is written, stroke by stroke}

\begin{quote}
\zh{北}
\end{quote}

\textbf{Two halves, left before right.}

\textbf{The left half, three strokes.} A vertical. A short bar from the left into it. Then a bar rising to the right into its foot, from the left.

\textbf{The right half, two strokes}, as on the right of \zh{比}: a stroke falling to the left, then a vertical that bends to the right along the foot and hooks up.

3 + 2 = five.

\subsection*{Your turn}
\begin{itemize}
\raggedright
  \item \emph{Write it:} \zh{北}, all five, with the model in front of you
  \item \emph{Cover:} write it again from memory, and count the pieces, not the strokes
\end{itemize}

\begin{tcolorbox}[breakable,colback=teal!4,colframe=teal!35!black,title={Before you move on}]
How many strokes in \zh{北}? (\textbf{Five}.) What is the right half? (\textbf{The right half of \zh{比}}.)

Source: \href{https://raw.githubusercontent.com/chanind/hanzi-writer-data/68d10a4b21150cae5e1ebbd223eed289cf32d90c/data/\%E5\%8C\%97.json}{Hanzi Writer Data \zh{北}, pinned stroke paths and medians}.
\end{tcolorbox}
\section[\texorpdfstring{\zh{西北} (xīběi)}{西北 (xīběi)}]{\zh{西北} (xīběi) — northwest; the Northwest (of China)}

\label{lesson:ZH-C422-xibei}



\begin{quote}
Before the new one: say the Mandarin for to navigate by satnav, then the Mandarin for the northeast.
\end{quote}

\subsection*{You'll want to know: \zh{西北}}
\textbf{\zh{西北}} — \emph{xīběi} — \textquotedblleft{}northwest; the Northwest (of China)\textquotedblright{}.

\begin{quote}
\textbf{\zh{风从西北来}} \emph{fēng cóng xīběi lái} — The wind is coming from the northwest.
\end{quote}

\begin{cousinweb}
\textbf{\zh{西北}} (\emph{xīběi}) is built from \zh{西} 'west' and \zh{北} 'north', both inherited from Middle Chinese, \emph{sej} and \emph{pok} (Guangyun, in Baxter's transcription).
\end{cousinweb}

\begin{grammarlens}[title={What you've built}]
One more word for getting around.
\end{grammarlens}

\subsection*{Your turn}
\begin{itemize}
\raggedright
  \item \emph{Say it:} \emph{xīběi}
  \item \emph{Say it:} \emph{xīběi}, once more
  \item \emph{Say it:} \emph{fēng cóng xīběi lái} — The wind is coming from the northwest.
  \item \emph{Recall:} say the Mandarin for to navigate by satnav, then the Mandarin for the northeast, then say \emph{xīběi} again
\end{itemize}

\begin{tcolorbox}[breakable,colback=teal!4,colframe=teal!35!black,title={Before you move on}]
What does \zh{西北} mean? (\textquotedblleft{}Northwest; the Northwest (of China)\textquotedblright{}.) Say it once more.
\end{tcolorbox}
\section[\texorpdfstring{\zh{北上} (běishàng)}{北上 (běishàng)}]{\zh{北上} (běishàng) — to go north; to head up north}

\label{lesson:ZH-C422-beishang}



\begin{quote}
Before the new one: say the Mandarin for the northeast, then the Mandarin for the northwest.
\end{quote}

\subsection*{You'll want to know: \zh{北上}}
\textbf{\zh{北上}} — \emph{běishàng} — \textquotedblleft{}to go north; to head up north\textquotedblright{}.

\begin{quote}
\textbf{\zh{明天我们坐飞机北上}} \emph{míngtiān wǒmen zuò fēijī běishàng} — Tomorrow we fly north.
\end{quote}

\begin{cousinweb}
\textbf{\zh{北上}} (\emph{běishàng}) is built from \zh{北} 'north' and \zh{上} 'to go up', both inherited from Middle Chinese, \emph{pok} and \emph{dzyangX} (Guangyun, in Baxter's transcription), which glosses the second 'to climb'.
\end{cousinweb}

\begin{grammarlens}[title={What you've built}]
One more word for getting around.
\end{grammarlens}

\subsection*{Your turn}
\begin{itemize}
\raggedright
  \item \emph{Say it:} \emph{běishàng}
  \item \emph{Say it:} \emph{běishàng}, once more
  \item \emph{Say it:} \emph{míngtiān wǒmen zuò fēijī běishàng} — Tomorrow we fly north.
  \item \emph{Recall:} say the Mandarin for the northeast, then the Mandarin for the northwest, then say \emph{běishàng} again
\end{itemize}

\begin{tcolorbox}[breakable,colback=teal!4,colframe=teal!35!black,title={Before you move on}]
What does \zh{北上} mean? (\textquotedblleft{}To go north; to head up north\textquotedblright{}.) Say it once more.
\end{tcolorbox}
\section[\texorpdfstring{\zh{快速} (kuàisù)}{快速 (kuàisù)}]{\zh{快速} (kuàisù) — rapid; high-speed}

\label{lesson:ZH-C422-kuaisu}



\begin{quote}
Before the new one: say the Mandarin for the northwest, then the Mandarin for to head up north.
\end{quote}

\subsection*{You'll want to know: \zh{快速}}
\textbf{\zh{快速}} — \emph{kuàisù} — \textquotedblleft{}rapid; high-speed\textquotedblright{}.

\begin{quote}
\textbf{\zh{我们坐快速公交去机场}} \emph{wǒmen zuò kuàisù gōngjiāo qù jīchǎng} — We take the rapid bus to the airport.
\end{quote}

\begin{cousinweb}
\textbf{\zh{快速}} (\emph{kuàisù}) is built from \zh{快} 'fast' and \zh{速} 'swift', both inherited from Middle Chinese, \emph{khwaejH} and \emph{suwk} (Guangyun, in Baxter's transcription).
\end{cousinweb}

\begin{grammarlens}[title={What you've built}]
One more word for getting around.
\end{grammarlens}

\subsection*{Your turn}
\begin{itemize}
\raggedright
  \item \emph{Say it:} \emph{kuàisù}
  \item \emph{Say it:} \emph{kuàisù}, once more
  \item \emph{Say it:} \emph{wǒmen zuò kuàisù gōngjiāo qù jīchǎng} — We take the rapid bus to the airport.
  \item \emph{Recall:} say the Mandarin for the northwest, then the Mandarin for to head up north, then say \emph{kuàisù} again
\end{itemize}

\begin{tcolorbox}[breakable,colback=teal!4,colframe=teal!35!black,title={Before you move on}]
What does \zh{快速} mean? (\textquotedblleft{}Rapid; high-speed\textquotedblright{}.) Say it once more.
\end{tcolorbox}
